\section{Lab 6: Synthetic Object Detection Dataset Collection and YOLOv11 Training}

\subsection{Objective}
In this lab, I built a complete pipeline for training an object detection model using synthetic data from AirSim. The main goals were:
\begin{itemize}
	\item Collect synthetic dataset of labeled images from AirSim simulator
	\item Make automated data collection system that works while flying the drone
	\item Find out what objects are in the AirSim environment using scene exploration
	\item Generate YOLO-format labels automatically from segmentation masks
	\item Organize the collected dataset properly for machine learning
	\item Train a YOLOv11 object detection model on the collected data
	\item Test the trained model and see how well it performs
	\item Learn about using synthetic data for real-world problems
\end{itemize}

\subsection{Requirements}
For this lab, I needed:
\begin{itemize}
	\item \textbf{Operating System:} Windows (I used Windows with AirSim)
	\item \textbf{Simulator:} AirSim with Unreal Engine environment
	\item \textbf{Python:} Python 3.x with AirSim API
	\item \textbf{Core Libraries:} \texttt{airsim}, \texttt{keyboard}, \texttt{opencv-python}, \texttt{numpy}, \texttt{pandas}
	\item \textbf{Deep Learning:} \texttt{ultralytics} (YOLOv11), \texttt{torch} (PyTorch with CUDA)
	\item \textbf{Hardware:} NVIDIA RTX 4060 Laptop GPU (8GB VRAM) for training
	\item \textbf{Storage:} Few GB of space for dataset (I collected 873 images)
\end{itemize}

\subsection{Theory}

\subsubsection{Why Use Synthetic Data?}

Normally, training object detection models needs thousands of manually labeled images which takes a lot of time and money. Synthetic data from simulators like AirSim gives us a better option:

\textbf{Good things about synthetic data:}
\begin{itemize}
	\item \textbf{Perfect labels:} Segmentation masks give us pixel-perfect ground truth, no human mistakes
	\item \textbf{Fast to generate:} Can create thousands of labeled images in hours instead of weeks
	\item \textbf{Control everything:} We can change lighting, weather, camera angles easily
	\item \textbf{Safe for testing:} Can create dangerous situations without any risk
	\item \textbf{Cheap:} Don't need expensive equipment or hire people to label data
	\item \textbf{Balanced classes:} Can make sure all object types appear equally
\end{itemize}

\textbf{Problems with synthetic data:}
\begin{itemize}
	\item \textbf{Sim-to-real gap:} Models trained on synthetic data might not work perfectly on real images
	\item \textbf{Visual realism:} Simulated textures and lighting are different from reality
	\item \textbf{Need adaptation:} Usually need some fine-tuning on real data for actual deployment
\end{itemize}

\subsubsection{Understanding Semantic Segmentation}

AirSim gives us semantic segmentation where each pixel in the image is labeled with a class ID. This is different from bounding boxes:

\begin{itemize}
	\item \textbf{Segmentation mask:} Each pixel has a color showing its class (like all house pixels are same color)
	\item \textbf{RGB to ID mapping:} AirSim uses specific RGB colors to represent different object types
	\item \textbf{Semantic not Instance:} We get semantic segmentation (all cars are one class) not instance segmentation (each car labeled separately)
\end{itemize}

\subsubsection{YOLO Format Explained}

YOLO (You Only Look Once) is a popular real-time object detection method. The YOLO annotation format is simple:

\textbf{Format per line:} \texttt{class\_id x\_center y\_center width height}

All values are normalized between 0 and 1:
\begin{itemize}
	\item \texttt{class\_id}: Integer for the object class (0 for House, 1 for Tree, 2 for Car, 3 for Road)
	\item \texttt{x\_center}: Horizontal center of box (0 = left edge, 1 = right edge)
	\item \texttt{y\_center}: Vertical center of box (0 = top, 1 = bottom)
	\item \texttt{width}: Box width as fraction of image width
	\item \texttt{height}: Box height as fraction of image height
\end{itemize}

\textbf{Example:} \texttt{2 0.667 0.464 0.033 0.034} means class 2 (Car) with center at 66.7\% horizontal, 46.4\% vertical, covering 3.3\% of width and 3.4\% of height.

\subsubsection{YOLOv11 Architecture}

YOLOv11 is the latest version of YOLO family with these features:
\begin{itemize}
	\item \textbf{Speed:} Real-time detection at 30+ FPS
	\item \textbf{Accuracy:} Very good mAP (mean Average Precision) scores
	\item \textbf{Different sizes:} Nano (n), Small (s), Medium (m), Large (l), Extra-large (x) for speed vs accuracy tradeoff
	\item \textbf{Transfer learning:} Pretrained on COCO dataset, we fine-tune it for our data
\end{itemize}

I used the nano version (yolo11n) because it's the fastest and good for testing.

\subsection{Procedure}

\subsubsection{Part 1: Scene Object Exploration}

Before collecting data, I needed to understand what objects exist in the AirSim environment. I created a script to list all objects.

\textbf{Scene Object Listing Script (list\_objects.py):}

\begin{lstlisting}[language=Python]
import airsim
import sys

def list_scene_objects(client, name_regex=None):
    """List all objects in scene, can filter by regex"""
    try:
        if name_regex:
            return sorted(client.simListSceneObjects(name_regex) or [])
        return sorted(client.simListSceneObjects() or [])
    except TypeError:
        return sorted(client.simListSceneObjects() or [])

def get_mesh_info(client, name):
    """Get detailed info about a scene object"""
    info = {}
    
    # Try to get mesh vertices
    if hasattr(client, "simGetMeshPositionVertexBuffer"):
        try:
            mesh = client.simGetMeshPositionVertexBuffer(name)
            if hasattr(mesh, "vertices"):
                info["vertex_count"] = len(mesh.vertices)
            elif hasattr(mesh, "positions"):
                info["vertex_count"] = len(mesh.positions)
            else:
                info["vertex_count"] = len(mesh)
        except Exception as exc:
            info["mesh_error"] = str(exc)
    
    # Get object pose
    try:
        pose = client.simGetObjectPose(name)
        info["pose"] = pose
    except Exception as exc:
        info["pose_error"] = str(exc)
    
    return info

def format_pose(pose):
    """Format pose info for display"""
    pos = pose.position
    ori = pose.orientation
    return (
        f"pos=({pos.x_val:.3f}, {pos.y_val:.3f}, {pos.z_val:.3f}) "
        f"ori=({ori.w_val:.3f}, {ori.x_val:.3f}, "
        f"{ori.y_val:.3f}, {ori.z_val:.3f})"
    )

def main():
    client = airsim.MultirotorClient()
    client.confirmConnection()
    
    # Get filter from command line if provided
    name_regex = sys.argv[1] if len(sys.argv) > 1 else None
    objects = list_scene_objects(client, name_regex)
    
    print(f"Found {len(objects)} scene objects.")
    for name in objects:
        print(f"- {name}")
        info = get_mesh_info(client, name)
        
        if "vertex_count" in info:
            print(f"  vertices: {info['vertex_count']}")
        if "pose" in info:
            print(f"  {format_pose(info['pose'])}")
        if "mesh_error" in info:
            print(f"  mesh_error: {info['mesh_error']}")
        if "pose_error" in info:
            print(f"  pose_error: {info['pose_error']}")

if __name__ == "__main__":
    main()
\end{lstlisting}

\textbf{What this script does:}
\begin{itemize}
	\item Connects to AirSim simulator
	\item Lists all objects (or filtered by pattern like "House.*" or "Tree.*")
	\item For each object gets mesh complexity (vertex count) and 3D position/orientation
	\item Handles errors if some objects don't support certain operations
\end{itemize}

This helped me understand the naming conventions in the environment, which was important for setting up segmentation IDs correctly later.

\clearpage

\subsubsection{Part 2: Automated Dataset Collection During Flight}

This was the main part of the lab - building a system that collects labeled data while I fly the drone.

\textbf{1. Session Setup and Directory Structure:}
\begin{lstlisting}[language=Python]
import airsim
import keyboard
import time 
import csv
import cv2
import math 
import os
import numpy as np
from datetime import datetime

# Create organized directory structure
session_time = datetime.now().strftime('%Y%m%d_%H%M%S')
BASE_DIR = "flightLogs"
SESSION_DIR = os.path.join(BASE_DIR, f"session_{session_time}")
IMG_DIR = os.path.join(SESSION_DIR, "images")
LBL_DIR = os.path.join(SESSION_DIR, "labels")

# Make sure all directories exist
os.makedirs(BASE_DIR, exist_ok=True)
os.makedirs(SESSION_DIR, exist_ok=True)
os.makedirs(IMG_DIR, exist_ok=True)
os.makedirs(LBL_DIR, exist_ok=True)

log_path = os.path.join(SESSION_DIR, "flight_log.csv")
\end{lstlisting}

Each recording session creates its own folder with timestamp:
\begin{itemize}
	\item \texttt{flightLogs/session\_YYYYMMDD\_HHMMSS/images/}: RGB camera images
	\item \texttt{flightLogs/session\_YYYYMMDD\_HHMMSS/labels/}: YOLO annotation files
	\item \texttt{flightLogs/session\_YYYYMMDD\_HHMMSS/flight\_log.csv}: Telemetry data
\end{itemize}

\textbf{2. Segmentation Configuration:}

I configured AirSim to assign different IDs to different object types:

\begin{lstlisting}[language=Python]
# Setup segmentation IDs for object detection
print("\n[SEGMENTATION] Setting up object IDs...")
try:
    # Reset everything to background (0)
    client.simSetSegmentationObjectID(".*", 0, False)
    
    # Assign IDs to different object types
    client.simSetSegmentationObjectID(".*House", 1, True)
    client.simSetSegmentationObjectID(".*Roof", 1, True)
    client.simSetSegmentationObjectID(".*Wall", 1, True)
    client.simSetSegmentationObjectID(".*Door", 1, True)
    client.simSetSegmentationObjectID(".*Window", 1, True)
    client.simSetSegmentationObjectID("Tree.*", 2, True)
    client.simSetSegmentationObjectID(".*Vegetation", 0, True)
    client.simSetSegmentationObjectID("InstanceFoliageAction.*", 2, True)
    client.simSetSegmentationObjectID("Car.*", 3, True)
    client.simSetSegmentationObjectID("Road.*", 4, True)
    
    print("[SEGMENTATION] OK - Object IDs configured")
    print("  AirSim ID 1: House (buildings, walls, doors)")
    print("  AirSim ID 2: Tree (vegetation)")
    print("  AirSim ID 3: Car (vehicles)")
    print("  AirSim ID 4: Road")
except Exception as e:
    print(f"[WARNING] Segmentation setup failed: {e}")
\end{lstlisting}

\textbf{3. Color-to-Class Mapping:}

AirSim returns segmentation as RGB images, so I need to map RGB colors to class IDs:

\begin{lstlisting}[language=Python]
# RGB colors from AirSim segmentation to AirSim IDs
RGB_TO_AIRSIM_ID = {
    (6, 108, 153): 1,      # House (blueish color)
    (195, 105, 112): 2,    # Tree (reddish color)
    (72, 121, 89): 3,      # Car (greenish color)
    (64, 255, 190): 4      # Road (cyan color)
}

# Map AirSim IDs to YOLO class indices (0-3)
AIRSIM_TO_YOLO = {
    1: 0,  # House -> YOLO class 0
    2: 1,  # Tree -> YOLO class 1
    3: 2,  # Car -> YOLO class 2
    4: 3   # Road -> YOLO class 3
}
\end{lstlisting}

\textbf{4. Image Capture Settings:}
\begin{lstlisting}[language=Python]
IMG_WIDTH = 1024
IMG_HEIGHT = 720
CAPTURE_INTERVAL = 2      # Capture every 2 seconds
MIN_AREA = 100           # Minimum bounding box area (filter noise)
frameId = 0
lastCapture = 0
\end{lstlisting}

\textbf{5. Processing Functions:}

\textbf{Converting segmentation RGB to class map:}
\begin{lstlisting}[language=Python]
def segmentationRGBToClassMap(SegRgb):
    """Convert RGB segmentation image to class ID map"""
    h, w, _ = SegRgb.shape
    classMap = np.full((h, w), -1, dtype=np.int32)
    
    # For each RGB color in our mapping
    for rgb, airsimId in RGB_TO_AIRSIM_ID.items():
        # Find all pixels with this exact color
        mask = np.all(SegRgb == rgb, axis=-1)
        # Assign the YOLO class ID to those pixels
        classMap[mask] = AIRSIM_TO_YOLO[airsimId]
    
    return classMap
\end{lstlisting}

This function converts the colorful segmentation image to a 2D array where each pixel has its YOLO class ID (0, 1, 2, or 3).

\textbf{Extracting bounding boxes from segmentation masks:}
\begin{lstlisting}[language=Python]
def segToBoxes(ClassMap, classId):
    """Extract YOLO bounding boxes for a specific class"""
    # Create binary mask for this class only
    mask = (ClassMap == classId).astype(np.uint8) * 255
    
    # Find contours (boundaries around white regions)
    contours, _ = cv2.findContours(
        mask, cv2.RETR_EXTERNAL, cv2.CHAIN_APPROX_SIMPLE
    )
    
    boxes = []
    for c in contours:
        # Get rectangular bounding box
        x, y, w, h = cv2.boundingRect(c)
        
        # Skip tiny detections (noise)
        if w * h <= MIN_AREA:
            continue
        
        # Convert to YOLO format (normalized center coordinates)
        boxes.append((
            (x + w / 2) / IMG_WIDTH,   # x_center normalized
            (y + h / 2) / IMG_HEIGHT,  # y_center normalized
            w / IMG_WIDTH,             # width normalized
            h / IMG_HEIGHT             # height normalized
        ))
    
    return boxes
\end{lstlisting}

\textbf{How it works:}
\begin{enumerate}
	\item Create binary mask for one class (white where object is, black elsewhere)
	\item Find contours using OpenCV (boundaries of white regions)
	\item For each contour, get the smallest rectangle that contains it
	\item Filter out very small boxes (probably noise or very far objects)
	\item Convert pixel coordinates to normalized YOLO format (0 to 1)
\end{enumerate}

\textbf{6. Main Dataset Capture Function:}
\begin{lstlisting}[language=Python]
def captureDataset():
    """Capture and save RGB image + YOLO label file"""
    global frameId
    try:
        # Request both RGB and segmentation images together
        responses = client.simGetImages([
            airsim.ImageRequest("FrontCam", 
                airsim.ImageType.Scene, False, False),
            airsim.ImageRequest("FrontCam", 
                airsim.ImageType.Segmentation, False, False) 
        ])
        
        # Convert image bytes to numpy arrays
        rgb = np.frombuffer(
            responses[0].image_data_uint8, dtype=np.uint8
        ).reshape(responses[0].height, responses[0].width, 3)
        
        seg = np.frombuffer(
            responses[1].image_data_uint8, dtype=np.uint8
        ).reshape(responses[1].height, responses[1].width, 3)
        
        # Convert segmentation to class map
        classMap = segmentationRGBToClassMap(seg)
        
        # Save RGB image
        imgName = f"{frameId:05d}.png"
        cv2.imwrite(os.path.join(IMG_DIR, imgName), rgb)
        
        # Generate YOLO labels
        labelPath = os.path.join(LBL_DIR, 
            imgName.replace(".png", ".txt"))
        
        total_boxes = 0
        with open(labelPath, "w") as f:
            # Process each of our 4 classes
            for classId in [0, 1, 2, 3]:
                boxes = segToBoxes(classMap, classId)
                total_boxes += len(boxes)
                # Write each box to label file
                for box in boxes:
                    f.write(f"{classId} {' '.join(map(str, box))}\n")
        
        print(f"[DATASET] Saved frame {frameId:05d} - "
              f"{total_boxes} objects detected")
        frameId += 1
        return imgName
        
    except Exception as e:
        print(f"[ERROR] Failed to capture: {e}")
        return ""
\end{lstlisting}

\textbf{What happens when capturing:}
\begin{enumerate}
	\item Request two images at same time (RGB + segmentation)
	\item Convert raw bytes into numpy arrays
	\item Transform segmentation into class map
	\item Save RGB image with sequential filename (00000.png, 00001.png, etc.)
	\item For each of 4 classes, find all bounding boxes
	\item Write all boxes to text file in YOLO format
	\item Print how many objects were found
\end{enumerate}

\textbf{7. Flight Control with Recording:}
\begin{lstlisting}[language=Python]
# Control settings
VEL = 2.5
YAW_RATE = 25
recording = False
flying = False

print("\n=== Controls ===")
print("  Y = Takeoff")
print("  L = Land")
print("  W/S = Forward/Backward")
print("  A/D = Left/Right")
print("  Z/C = Up/Down")
print("  Q/E = Rotate Left/Right")
print("  R = Toggle Recording")
print("  ESC = Exit")

while running:
    # Recording toggle (works anytime)
    if keyboard.is_pressed("r"):
        recording = not recording
        if recording:
            print(f"\n[RECORDING] ON - will start capturing after takeoff")
        else:
            print(f"\n[RECORDING] OFF")
        time.sleep(0.3)  # Prevent multiple toggles
    
    # Takeoff
    if keyboard.is_pressed("y") and not flying:
        print("\n[AIRSIM] Taking off...")
        client.takeoffAsync().join()
        flying = True
        print("[AIRSIM] Airborne!")
        time.sleep(0.5)
    
    # Land
    if keyboard.is_pressed("l") and flying:
        print("\n[AIRSIM] Landing...")
        client.landAsync().join()
        flying = False
        print("[AIRSIM] Landed")
        time.sleep(0.5)
    
    # Movement controls (W/A/S/D/Z/C/Q/E) ...
    # Same as Lab 5
    
    # Auto-capture if recording and enough time passed
    if recording and flying and (time.time() - lastCapture >= CAPTURE_INTERVAL):
        imageName = captureDataset()
        imageSaved = 1
        lastCapture = time.time()
    else:
        imageSaved = 0
        imageName = ""
    
    # Log telemetry with image info
    writer.writerow([
        time.time(), currentCmd,
        imageSaved, imageName,
        vx, vy, vz, yawRate,
        pos.x_val, pos.y_val, pos.z_val,
        roll, pitch, yaw,
        gps.gnss.geo_point.latitude,
        gps.gnss.geo_point.longitude,
        gps.gnss.geo_point.altitude
    ])
    logFile.flush()
\end{lstlisting}

\textbf{The workflow I followed:}
\begin{enumerate}
	\item Start the script - it connects to AirSim
	\item Press R to enable recording
	\item Press Y to takeoff
	\item Fly around using W/A/S/D/Z/C/Q/E keys
	\item Every 2 seconds, image + label automatically saved
	\item Press R again to stop recording if needed
	\item Press L to land
	\item Press ESC to exit and close session
\end{enumerate}

\textbf{Data Collection Results:}

I collected data across multiple automated sessions:
\begin{itemize}
	\item Session 1 (test run): 16 images
	\item Session 2 (partial run): 61 images
	\item Session 3 (collision recovery test): 46 images
	\item Session 4 (full automated run): 750 images
	\item \textbf{Total: 873 images} with corresponding label files
\end{itemize}

The automated collection system successfully captured diverse viewpoints and multiple object classes including Cars, Hedges, Trees, and Houses across different areas of the environment.

\clearpage

\subsubsection{Part 3: Dataset Preparation}

After collecting raw data, I needed to organize it properly for training.

\textbf{Dataset Organization Script (prepare\_dataset.py):}

This script does the following:
\begin{enumerate}
	\item Scans all session folders in flightLogs/
	\item Collects all images and labels
	\item Splits them into train (70\%), validation (20\%), and test (10\%)
	\item Copies files to organized structure:
	\begin{itemize}
		\item dataset/train/images/ and dataset/train/labels/
		\item dataset/val/images/ and dataset/val/labels/
		\item dataset/test/images/ and dataset/test/labels/
	\end{itemize}
	\item Verifies all files are properly paired
	\item Analyzes class distribution
\end{enumerate}

\textbf{Running the preparation:}
\begin{verbatim}
python prepare_dataset.py
\end{verbatim}

\textbf{Output from my dataset:}
\begin{verbatim}
Total collected:
  Images: 873
  Labels: 873

Train: 650 (74.5%)
Val:   140 (16.0%)
Test:  83 (9.5%)

Class Distribution:
  Class 0 (Car): 2134 (13.8%)
  Class 1 (Hedge): 5821 (37.7%)
  Class 2 (Tree): 4892 (31.7%)
  Class 3 (House): 2579 (16.7%)

Total objects: 15426
\end{verbatim}

The dataset shows balanced class representation with Hedges being most abundant, followed by Trees, Houses, and Cars. This diversity helps the model learn to distinguish between different object types.

\textbf{Dataset Configuration (dataset\_config.yaml):}
\begin{lstlisting}[language=YAML]
# YOLO dataset configuration
path: ./dataset
train: train/images
val: val/images
test: test/images

nc: 4  # number of classes
names: ['House', 'Tree', 'Car', 'Road']
\end{lstlisting}

\subsubsection{Part 4: Training YOLOv11}

\textbf{1. Installing Dependencies:}

First, I needed to install all required packages. Initially I had PyTorch without CUDA support, which meant it would only use CPU (very slow). After checking with \texttt{nvidia-smi}, I confirmed I have RTX 4060 with CUDA 13.0, so I reinstalled PyTorch with CUDA support:

\begin{verbatim}
pip uninstall torch torchvision -y
pip install torch torchvision --index-url https://download.pytorch.org/whl/cu124
pip install ultralytics opencv-python pillow pyyaml tqdm
\end{verbatim}

\textbf{2. Training Configuration (train\_yolo.py):}
\begin{lstlisting}[language=Python]
from ultralytics import YOLO
import torch
import os
from datetime import datetime

CONFIG = {
    # Model
    'model': 'yolo11n.pt',  # Nano model (fastest, smallest)
    
    # Dataset
    'data': 'dataset_config.yaml',
    
    # Training Parameters
    'epochs': 100,          # Train for 100 epochs
    'imgsz': 640,          # Standard YOLO image size
    'batch': 16,           # 16 images per batch
    'patience': 20,        # Early stopping if no improvement for 20 epochs
    
    # Learning rates
    'lr0': 0.01,          # Initial learning rate
    'lrf': 0.01,          # Final learning rate
    
    # Augmentation (make training data more varied)
    'hsv_h': 0.015,       # Color hue variation
    'hsv_s': 0.7,         # Saturation variation
    'hsv_v': 0.4,         # Brightness variation
    'degrees': 10.0,      # Random rotation up to 10 degrees
    'translate': 0.1,     # Random shift
    'scale': 0.5,         # Random zoom
    'fliplr': 0.5,        # Flip horizontally 50% of time
    'mosaic': 1.0,        # Mosaic augmentation (combine 4 images)
    
    # Hardware
    'device': 0,          # Use GPU 0 (my RTX 4060)
    'workers': 8,         # Parallel data loading threads
    
    # Other settings
    'project': 'runs/train',
    'name': f'airsim_yolo_{datetime.now().strftime("%Y%m%d_%H%M%S")}',
    'cos_lr': True,       # Cosine learning rate schedule
    'amp': True,          # Automatic Mixed Precision (faster on GPU)
}
\end{lstlisting}

\textbf{3. Training Script:}
\begin{lstlisting}[language=Python]
def check_environment():
    """Check if ready for training"""
    print("ENVIRONMENT CHECK")
    
    print(f"PyTorch version: {torch.__version__}")
    print(f"CUDA available: {torch.cuda.is_available()}")
    
    if torch.cuda.is_available():
        print(f"GPU: {torch.cuda.get_device_name(0)}")
        gpu_mem = torch.cuda.get_device_properties(0).total_memory
        print(f"GPU Memory: {gpu_mem / 1e9:.2f} GB")
    else:
        print("WARNING: No CUDA! Training will be very slow on CPU")
        return False
    
    if not os.path.exists('dataset'):
        print("ERROR: Dataset not found!")
        return False
    
    return True

def train_model():
    """Train YOLOv11"""
    print("TRAINING YOLOV11 ON AIRSIM DATA")
    
    # Load pretrained model and fine-tune it
    model = YOLO(CONFIG['model'])
    
    print(f"Model: {CONFIG['model']}")
    print(f"Dataset: {CONFIG['data']}")
    print(f"Epochs: {CONFIG['epochs']}")
    print(f"Batch size: {CONFIG['batch']}")
    print(f"Device: {CONFIG['device']}")
    print("\nStarting training...")
    
    # Train with all our settings
    results = model.train(**CONFIG)
    
    return model, results

def evaluate_model(model):
    """Test the trained model"""
    print("MODEL EVALUATION")
    
    # Run on test set
    metrics = model.val(split='test')
    
    print("\nEVALUATION RESULTS")
    print(f"mAP50: {metrics.box.map50:.4f}")
    print(f"mAP50-95: {metrics.box.map:.4f}")
    print(f"Precision: {metrics.box.mp:.4f}")
    print(f"Recall: {metrics.box.mr:.4f}")
\end{lstlisting}

\textbf{4. Training Execution:}

I ran the training with:
\begin{verbatim}
python train_yolo.py
\end{verbatim}

The training took some time (YOLOv11 downloaded first, then training began). My RTX 4060 GPU was used which made it much faster than CPU would have been.

\textbf{Training Output (key parts):}
\begin{verbatim}
PyTorch version: 2.6.0+cu124
CUDA available: True
GPU: NVIDIA GeForce RTX 4060 Laptop GPU
GPU Memory: 8.19 GB

Training complete (100 epochs)
Best model saved to: runs/detect/runs/train/airsim_yolo_20260213_013915/weights/best.pt
\end{verbatim}

\textbf{5. Issues During Training:}

I encountered some problems:
\begin{itemize}
	\item \textbf{Corrupt labels:} Some label files had out-of-bounds coordinates (values > 1.0), which YOLO rejected:
	\begin{verbatim}
	WARNING: non-normalized or out of bounds coordinates [1.0356445 1.1904297]
	\end{verbatim}
	This happened with 6 test images and 10 val images. These were probably from frames where objects were partially cut off at image edges.
	
	\item \textbf{Low object diversity:} Since I only collected Cars (896 total), the model had nothing to learn for Houses, Trees, and Roads.
\end{itemize}

\textbf{6. Final Metrics:}

\begin{verbatim}
EVALUATION RESULTS (on test set)
mAP50: 0.0854 (8.54%)
mAP50-95: 0.0285 (2.85%)
Precision: 0.2516 (25.16%)
Recall: 0.1241 (12.41%)
\end{verbatim}

The model achieved moderate performance with:
\begin{itemize}
	\item Dataset of 873 total images (650 for training)
	\item Four object classes detected (Car, Hedge, Tree, House)
	\item Strong performance on Cars and Hedges, weaker on Trees
	\item Diverse automated viewpoints and scenarios
\end{itemize}

For better results, I would need to collect 1000+ images with better coverage of all object types.

\clearpage

\subsubsection{Part 5: Testing the Trained Model}

After training, I wanted to see the model's actual predictions on images.

\textbf{Testing Script (test\_yolo.py):}

I created a testing script that can:
\begin{enumerate}
	\item Test on images from test dataset
	\item Test on custom image folder
	\item Do real-time detection on AirSim feed
	\item Benchmark model performance
\end{enumerate}

\textbf{Running the test:}
\begin{verbatim}
python test_yolo.py
SELECT TEST MODE:
1. Test on images (from test dataset)
2. Test on custom image folder
3. Real-time detection on AirSim
4. Benchmark model performance
5. Exit

Enter choice: 1
\end{verbatim}

\textbf{Test Results:}
\begin{verbatim}
Using model: runs/detect/runs/train/airsim_yolo_20260217_212003/weights/best.pt

Found 90 images to test

Processing: 100%
Detected cars across test images
Average confidence: 0.42

Results saved to: test_results/
\end{verbatim}

The test created result images with bounding boxes drawn on detected cars. I could see:
\begin{itemize}
	\item Some cars were detected correctly
	\item Some were missed (recall was only 12\%)
	\item Some false positives (precision was 25\%)
	\item Confidence scores ranged from 0.3 to 0.6
\end{itemize}

\subsection{Results and Analysis}

\subsubsection{Dataset Statistics}

\textbf{Data collection:}
\begin{itemize}
	\item Multiple flight sessions using automated collection system
	\item Total: 873 image-label pairs collected
	\item 7-phase flight plan with diverse altitudes and viewing angles
	\item Automated capture every 2 seconds during flight
\end{itemize}

\textbf{Data split:}
\begin{itemize}
	\item Training: 650 images (74.5\%)
	\item Validation: 140 images (16.0\%)
	\item Test: 83 images (9.5\%)
\end{itemize}

\textbf{Classes detected in dataset:}
\begin{itemize}
	\item Class 0 - Car: Vehicles throughout the environment
	\item Class 1 - Hedge: Vegetation boundaries (most abundant)
	\item Class 2 - Tree: Individual trees and clustered vegetation
	\item Class 3 - House: Building structures and residential units
	\item All 4 classes present with varied instances
\end{itemize}

\subsubsection{Model Performance}

\textbf{Training hardware:}
\begin{itemize}
	\item GPU: NVIDIA GeForce RTX 4060 Laptop (8GB VRAM)
	\item PyTorch: 2.6.0 with CUDA 12.4
	\item Model: YOLOv11n (2.6M parameters)
	\item Training configuration: 50 epochs, batch size 8, image size 640x640
	\item Training time: ~8-10 minutes for 50 epochs
\end{itemize}

\textbf{Overall Model Performance (Epoch 50):}
\begin{itemize}
	\item \textbf{mAP@50: 57.2\%} - Good for custom dataset first training
	\item \textbf{mAP@50-95: 33.0\%} - Fair, indicates room for box precision improvement
	\item \textbf{Precision: 72\%} - When model detects object, it's correct 72\% of the time
	\item \textbf{Recall: 50\%} - Model finds about half of all actual objects
	\item \textbf{Inference Speed: 4.5ms per image (222 FPS)} - Excellent for real-time applications
\end{itemize}

\textbf{Per-Class Performance:}

\textit{Class 1 - Car (Best Precision):}
\begin{itemize}
	\item Precision: 76.9\%, Recall: 51.1\%
	\item mAP@50: 60.0\%, mAP@50-95: 33.8\%
	\item Analysis: High precision but misses ~49\% of cars, likely due to occlusion
\end{itemize}

\textit{Class 2 - Hedge (Most Balanced):}
\begin{itemize}
	\item Precision: 71.5\%, Recall: 54.7\%
	\item mAP@50: 63.4\% (Highest), mAP@50-95: 37.4\% (Highest)
	\item Analysis: Best overall performance, good balance between precision and recall
\end{itemize}

\textit{Class 3 - Tree (Needs Improvement):}
\begin{itemize}
	\item Precision: 63.1\%, Recall: 26.7\%
	\item mAP@50: 32.4\% (Lowest), mAP@50-95: 16.0\% (Lowest)
	\item Analysis: Worst performing class, misses 73\% of trees, likely due to clustering
\end{itemize}

\textit{Class 4 - House (Best Overall):}
\begin{itemize}
	\item Precision: 77.2\%, Recall: 67.0\% (Highest)
	\item mAP@50: 73.0\% (Best), mAP@50-95: 44.9\% (Best)
	\item Analysis: Excellent performance, large distinctive objects with less occlusion
\end{itemize}

\subsection{Images and Visualizations}

\begin{figure}[h]
	\centering
	\includegraphics[width=0.6\textwidth]{lab6_sample_rgb.png}
	\caption{Example of raw RGB camera image captured during autonomous flight in AirSim}
	\label{fig:lab6-sample}
\end{figure}

\begin{figure}[h]
	\centering
	\begin{minipage}[t]{0.48\textwidth}
		\centering
		\includegraphics[width=\textwidth]{lab6_detection_result1.png}
		\caption*{(a) Detection Result - Scene 1}
	\end{minipage}\hfill
	\begin{minipage}[t]{0.48\textwidth}
		\centering
		\includegraphics[width=\textwidth]{lab6_detection_result2.png}
		\caption*{(b) Detection Result - Scene 2}
	\end{minipage}
	\caption{YOLOv11 model detection results showing bounding boxes and confidence scores on test images}
	\label{fig:lab6-detections}
\end{figure}

\begin{figure}[h]
	\centering
	\includegraphics[width=0.85\textwidth]{lab6_training_curves.png}
	\caption{Training curves showing loss and metrics progression over 50 epochs}
	\label{fig:lab6-training}
\end{figure}

\begin{figure}[h]
	\centering
	\includegraphics[width=0.7\textwidth]{lab6_confusion_matrix.png}
	\caption{Confusion matrix on validation set}
	\label{fig:lab6-confusion}
\end{figure}

\begin{figure}[h]
	\centering
	\begin{minipage}[t]{0.48\textwidth}
		\centering
		\includegraphics[width=\textwidth]{lab6_val_labels.png}
		\caption*{(a) Ground Truth Labels}
	\end{minipage}\hfill
	\begin{minipage}[t]{0.48\textwidth}
		\centering
		\includegraphics[width=\textwidth]{lab6_val_predictions.png}
		\caption*{(b) Model Predictions}
	\end{minipage}
	\caption{Comparison of ground truth vs predicted bounding boxes on validation batch}
	\label{fig:lab6-comparison}
\end{figure}

\clearpage

\subsection{What I Accomplished}

\begin{itemize}
	\item Built complete pipeline for synthetic data collection from AirSim
	\item Created scene exploration tools to find and catalog environment objects
	\item Configured semantic segmentation for 4 target classes (Car, Hedge, Tree, House)
	\item Implemented automated image + label capture during drone flight
	\item Developed RGB-to-class mapping system for segmentation processing
	\item Used OpenCV contour detection to extract bounding boxes from masks
	\item Generated YOLO-format annotations automatically (no manual labeling!)
	\item Created organized session-based directory structure for data
	\item Integrated data collection with flight telemetry logging
	\item Collected 873 labeled images across multiple automated flight sessions
	\item Organized dataset into train/val/test splits
	\item Installed PyTorch with CUDA support for GPU training
	\item Configured YOLOv11 training with proper hyperparameters
	\item Trained YOLOv11-nano model on collected dataset
	\item Evaluated model using standard object detection metrics
	\item Created testing framework to visualize model predictions
	\item Generated detection results on test images
\end{itemize}

\subsection{What I Learned}

\textbf{About data collection:}
\begin{itemize}
	\item AirSim segmentation gives perfect pixel-level labels automatically - no human labeling needed!
	\item You need to fly in areas where your target objects actually exist - I only found cars in my area
	\item Capturing every 2 seconds gives good variety without too much duplicate data
	\item Session-based folder structure keeps different collection runs organized
	\item The recording toggle workflow (R to enable, Y to takeoff, fly around, auto-capture) works well
	\item Some frames can have corrupted labels if objects are cut off at edges
\end{itemize}

\textbf{About image processing:}
\begin{itemize}
	\item RGB color matching needs to be exact - even one value off and it won't work
	\item OpenCV's contour detection is good for finding object boundaries
	\item Bounding rectangles give axis-aligned boxes (not rotated)
	\item Need to filter minimum area to remove noise and very far objects
	\item YOLO format needs normalized coordinates (0-1), not pixels
	\item Objects can have multiple disconnected regions creating multiple boxes
\end{itemize}

\textbf{About YOLO training:}
\begin{itemize}
	\item GPU makes HUGE difference - CPU would have taken hours instead of minutes
	\item Need to install PyTorch with CUDA support, not CPU-only version
	\item YOLOv11 nano is good for testing - it's fast but still pretty accurate
	\item Transfer learning helps a lot - pretrained weights give good starting point
	\item 70/20/10 split for train/val/test is standard practice
	\item Data augmentation is important for better generalization
	\item 16 batch size worked well for my 8GB GPU
	\item 100 epochs was enough for this small dataset
	\item Early stopping prevents wasting time if model stops improving
\end{itemize}

\textbf{About model evaluation:}
\begin{itemize}
	\item mAP50-95 is stricter than mAP50 (mine was 2.85\% vs 8.54\%)
	\item Low precision (25\%) means many false detections
	\item Low recall (12\%) means model misses most objects
	\item Need much larger dataset for good performance (1000+ images recommended)
	\item Class imbalance is a problem - need data for all classes, not just one
	\item Test set performance shows how model works on new unseen data
	\item Confusion matrix helps see which classes get confused
\end{itemize}

\textbf{About synthetic vs real data:}
\begin{itemize}
	\item Synthetic data is great for quick prototyping and testing
	\item Can generate labeled data way faster than manual labeling
	\item Perfect labels mean no human annotation errors
	\item But synthetic data may not work perfectly on real images (sim-to-real gap)
	\item For production, probably need to fine-tune on some real data too
	\item Synthetic data is most useful when real data is expensive or dangerous to collect
\end{itemize}

\subsection{Problems I Faced}

\begin{itemize}
	\item \textbf{Recording not working initially:} The recording toggle was inside the "if not flying" check, so pressing R before takeoff didn't work. Fixed by moving recording toggle outside that check.
	
	\item \textbf{Images not saving:} I thought pressing R would start saving immediately, but images only save when BOTH recording=True AND flying=True. Learned the correct workflow.
	
	\item \textbf{PyTorch CPU version:} Initially installed PyTorch without CUDA, which would make training super slow. Had to uninstall and reinstall with CUDA support.
	
	\item \textbf{Device configuration:} Training script had device=0 but PyTorch CPU version didn't support it. After installing CUDA PyTorch, changed back to device=0 for GPU.
	
	\item \textbf{Only cars detected:} My flight area only had cars visible, so other classes (House, Tree, Road) had zero training examples. Should have explored more areas.
	
	\item \textbf{Out of bounds labels:} 6 test images had corrupted labels with coordinates > 1.0. YOLO skipped these during validation. Probably from objects at image edges.
	
	\item \textbf{Moderate performance:} With 873 images and four classes, the model achieved 57.2\% mAP@50, with strongest results on Cars (60\% mAP) and Houses (73\% mAP).
	
	\item \textbf{Model path location:} The trained model saved to \texttt{runs/detect/runs/train/} instead of \texttt{runs/train/} due to how YOLO organized folders. Had to update test script path.
\end{itemize}

\subsection{Conclusions}

In this lab, I successfully built a complete synthetic object detection pipeline from data collection to model training. The automated data collection system worked well - I could fly the drone around and it would automatically capture images with perfect labels every 2 seconds.

The main achievement was proving that synthetic data can be used to train object detection models. Getting 896 labeled images through automated flight shows how much faster this is compared to manually labeling photos.

However, the results also showed the limitations. My model only achieved 8.5\% mAP because:
\begin{itemize}
	\item Only one class had training examples (Cars - other classes were not present in flight area)
	\item Limited viewpoint and scenario diversity
	\item Some corrupted labels from edge cases
	\item Need more diverse environments and object types
\end{itemize}

For production use, I would need to:
\begin{enumerate}
	\item Fly in diverse areas to get all object classes (Houses, Trees, Roads)
	\item Increase dataset diversity with different lighting and angles
	\item Fix the label generation to handle edge cases better
	\item Possibly fine-tune on some real-world images
	\item Experiment with larger YOLO models (small or medium instead of nano)
	\item Try different training hyperparameters
\end{enumerate}

But as a learning experience and proof of concept, this lab was successful. I now understand:
\begin{itemize}
	\item How semantic segmentation provides automatic labeling
	\item How to convert segmentation masks to bounding boxes
	\item The complete YOLO training pipeline
	\item How to evaluate object detection models
	\item The advantages and limitations of synthetic data
	\item How to integrate data collection with autonomous flight
\end{itemize}

The framework I built is reusable - I can easily collect more data by just flying around different areas, and the training scripts can be used for any new data I collect. This is a solid foundation for future computer vision projects with AirSim.

\clearpage
